\section{Study 3: Noise Models and Zero-Mean Invariance}
\label{sec:study3}

\subsection{Motivation}
\label{sec:study3:motivation}

The NumPy ballistic simulator is deterministic: given a commanded velocity
$\mathbf{v}_{\mathrm{cmd}}$, the landing position is a fixed function of that input.
A closed-form ballistic inverter could therefore solve the task in one step,
rendering the GP-based learning loop superfluous.
In a physical system, the gap between commanded and delivered velocity arises from
several sources: joint-torque saturation, gripper compliance, and release timing
uncertainty.
This study introduces explicit stochastic noise at the velocity-override interface to
make the RL problem non-trivial, characterises which noise types are learnable, and
proves a formal theorem on the boundary between learnable and unlearnable noise.

\subsection{Four Noise Models}
\label{sec:study3:models}

All noise models are realised through the \texttt{perturb\_numpy} interface:
\begin{equation}
  \mathbf{v}_{\mathrm{actual}} = \alpha\,\mathbf{v}_{\mathrm{cmd}} + \boldsymbol{\delta},
  \label{eq:perturb_interface}
\end{equation}
where the scale $\alpha$ and offset $\boldsymbol{\delta}$ are detached constants
(no gradient tape), so the policy gradient flows through $\mathbf{v}_{\mathrm{cmd}}$
unobstructed.

\paragraph{VelocityBiasNoise (Gaussian).}
Zero-mean additive perturbation:
\begin{equation}
  \alpha = 1, \quad \boldsymbol{\delta} \sim \mathcal{N}(\mathbf{0},\, \sigma^2 I_3).
  \label{eq:bias_noise}
\end{equation}
The expected delivered velocity equals the commanded velocity.
By Theorem~\ref{thm:zeromean} below, this noise type is unlearnable: the optimal
policy under Gaussian noise coincides with the noiseless optimum.

\paragraph{VelocitySlipNoise.}
Multiplicative under-delivery with additive scatter:
\begin{equation}
  \alpha = 1 - \alpha_{\mathrm{slip}}, \quad
  \boldsymbol{\delta} \sim \mathcal{N}(\mathbf{0},\, \sigma^2 I_3),
  \label{eq:slip_noise}
\end{equation}
where $\alpha_{\mathrm{slip}} \in (0, 1)$ is the slip fraction.
The mean delivered speed is $(1 - \alpha_{\mathrm{slip}})\,\|\mathbf{v}_{\mathrm{cmd}}\|$,
introducing a systematic undershoot that a noise-aware policy can compensate by
commanding $v_{\mathrm{cmd}} = v^* / (1 - \alpha_{\mathrm{slip}})$.

\paragraph{Salt-and-Pepper Noise.}
With probability $p$, a velocity component is replaced by a random spike of magnitude
$s$; otherwise it is delivered exactly.
This models electromagnetic glitches or encoder errors.
Unlike slip, salt-and-pepper does not have a constant mean shift across all
commanded speeds, so the compensation strategy is non-trivial.

\paragraph{ReleaseTimingJitter (paper-faithful).}
Gripper delay $t_d \sim \mathcal{U}(a, b)$ with $a > 0$ causes systematic undershoot:
\begin{equation}
  \mathbf{v}_{\mathrm{actual}} \approx
    \left(1 - \frac{\gamma\,t_d}{\|\mathbf{v}_{\mathrm{cmd}}\|}\right)
    \mathbf{v}_{\mathrm{cmd}},
  \label{eq:jitter_noise}
\end{equation}
matching the gripper-delay model of \cite{turcato2025mcpilot}.
This is a biased noise type ($\mathbb{E}[\alpha] < 1$) and hence learnable.

\subsection{Zero-Mean Noise Invariance Theorem}
\label{sec:study3:theory}

\begin{theorem}[Optimal Action Invariance Under Symmetric Noise]
\label{thm:zeromean}
Let $f: \mathbb{R}^3 \to \mathbb{R}^2$ be the deterministic ballistic map from release
velocity to landing position, smooth and invertible near the optimum.
Let $\boldsymbol{\epsilon} \sim p_\epsilon$ with $\mathbb{E}[\boldsymbol{\epsilon}] = \mathbf{0}$
and $p_\epsilon$ symmetric around zero.
Define the noise-aware objective
$J_{\mathrm{aware}}(\theta) = \mathbb{E}_{\boldsymbol{\epsilon}}[c(f(\mathbf{v}_{\mathrm{cmd}}
+ \boldsymbol{\epsilon}), \mathbf{P})]$
and the noiseless objective
$J_{\mathrm{naive}}(\theta) = c(f(\mathbf{v}_{\mathrm{cmd}}), \mathbf{P})$.
Then the minimisers of $J_{\mathrm{aware}}$ and $J_{\mathrm{naive}}$ coincide.
That is, symmetric zero-mean noise does not shift the optimal mean action of the policy.
\end{theorem}

\begin{proof}
By first-order Taylor expansion of $f$ around the optimal velocity $\mathbf{v}^*$:
\begin{equation*}
  f(\mathbf{v}^* + \boldsymbol{\epsilon}) \approx f(\mathbf{v}^*) + J_f \boldsymbol{\epsilon},
  \quad J_f = \tfrac{\partial f}{\partial \mathbf{v}}\big|_{\mathbf{v}^*}.
\end{equation*}
Taking the expectation:
\begin{equation*}
  \mathbb{E}_\epsilon[f(\mathbf{v}^* + \boldsymbol{\epsilon})]
  = f(\mathbf{v}^*) + J_f \underbrace{\mathbb{E}[\boldsymbol{\epsilon}]}_{=\,\mathbf{0}}
  = f(\mathbf{v}^*).
\end{equation*}
The gradient of the aware objective with respect to $\mathbf{v}_{\mathrm{cmd}}$ is:
\begin{align*}
  \nabla_{\mathbf{v}} J_{\mathrm{aware}}
  &= \nabla_{\mathbf{v}} \mathbb{E}_\epsilon[c(f(\mathbf{v} + \boldsymbol{\epsilon}))] \\
  &= \mathbb{E}_\epsilon\!\left[J_f^\top \nabla_{\mathbf{p}} c\bigl(f(\mathbf{v}
     + \boldsymbol{\epsilon})\bigr)\right].
\end{align*}
At the noiseless optimum $\mathbf{v}^*$ where $f(\mathbf{v}^*) = \mathbf{P}$, the
landing error $f(\mathbf{v}^* + \boldsymbol{\epsilon}) - \mathbf{P} \approx
J_f \boldsymbol{\epsilon}$ is an odd function of $\boldsymbol{\epsilon}$ to first order.
The cost gradient $\nabla_{\mathbf{p}} c(\mathbf{p})$ is itself an odd function of
the landing error $\mathbf{p} - \mathbf{P}$, so the composition
$\nabla_{\mathbf{p}} c(f(\mathbf{v}^* + \boldsymbol{\epsilon}))$ is odd in
$\boldsymbol{\epsilon}$ \emph{at} $\mathbf{v}^*$ (this argument applies specifically at
the optimum, not in general).
For symmetric $p_\epsilon$, the expectation of any odd function is zero:
$\mathbb{E}_\epsilon[\nabla_{\mathbf{p}} c(f(\mathbf{v}^* +
\boldsymbol{\epsilon}))] = \mathbf{0}$, which is the same stationarity condition as the
noiseless problem $\nabla_{\mathbf{v}} J_{\mathrm{naive}}\big|_{\mathbf{v}^*} = 0$.
Hence both objectives share the same stationary point $\mathbf{v}^*$.
\end{proof}

\begin{remark}
The result relies on the linearity of $f$ near $\mathbf{v}^*$ and the zero-mean
symmetry of $p_\epsilon$.
For skewed or heavy-tailed noise the approximation breaks and a genuine gap between
aware and naive policies can emerge.
The biased noise types (Slip, Jitter) fall outside this theorem because
$\mathbb{E}[\boldsymbol{\epsilon}] \neq \mathbf{0}$; they are learnable.
\end{remark}

\paragraph{GP regularisation consequence.}
An unexpected corollary of zero-mean noise appears in GP hyperparameter estimation.
With $\sigma = 0$ (zero noise), $N_{\mathrm{exp}} = 10$ perfectly deterministic
exploration throws give the GP log-likelihood a degenerate ridge: the MLE drives
lengthscales to $1791$\,m, $1602$\,m, and $357$\,m — three to four orders of magnitude
larger than the $\sim$1\,m domain.
A GP with these lengthscales predicts the same landing position regardless of launch
speed, making all policy gradients zero.
With $\sigma = 0.04$\,m/s, repeated throws at similar speeds produce slightly different
landings, and MLE finds finite lengthscales that explain the variance as a function of
speed.
\emph{Small noise is a necessary GP regulariser} for deterministic simulation data; we
recommend $\sigma_n \approx 0.04$\,m/s for the throwing domain.

\subsection{Empirical Comparison: Noise-Aware vs.\ Naive}
\label{sec:study3:empirical}

Table~\ref{tab:noise_results} reports hit rate, mean error, and final cost for all
tested noise conditions.
Aware and naive policies share identical exploration throws; the only difference is
whether noise is propagated through MC rollouts during policy optimisation.

\begin{table*}[t]
  \centering
  \caption{Performance under different arm-side noise models.
           Aware = noise propagated through MC rollouts; Naive = noiseless rollouts.
           Hit radius = 10\,cm. All results at seed=1, 4 policy trials.}
  \label{tab:noise_results}
  \small
  \begin{tabular}{llcccc}
    \toprule
    Noise Type & Level & Aware? & Hit Rate (\%) & Mean Error (m) & Final Cost \\
    \midrule
    Slip & $\alpha{=}0.05,\,\sigma{=}0.04$ & No  & 100 & 0.0682 & 0.00349 \\
    Slip & $\alpha{=}0.05,\,\sigma{=}0.04$ & Yes & 100 & 0.0260 & 0.00700 \\
    Slip & $\alpha{=}0.10,\,\sigma{=}0.04$ & No  & \textbf{25}  & 0.1242 & 0.00355 \\
    Slip & $\alpha{=}0.10,\,\sigma{=}0.04$ & Yes & \textbf{100} & 0.0450 & 0.00672 \\
    Slip & $\alpha{=}0.15,\,\sigma{=}0.04$ & No  & 100 & 0.0578 & 0.00339 \\
    Slip & $\alpha{=}0.15,\,\sigma{=}0.04$ & Yes & 100 & 0.0359 & 0.00623 \\
    Slip & $\alpha{=}0.20,\,\sigma{=}0.04$ & No  & \textbf{0}   & 0.1881 & 0.00404 \\
    Slip & $\alpha{=}0.20,\,\sigma{=}0.04$ & Yes & \textbf{100} & 0.0219 & 0.01097 \\
    \midrule
    Salt-Pepper & $p{=}0.05,\,s{=}0.30$ & No  & 75  & 0.0439 & 0.00410 \\
    Salt-Pepper & $p{=}0.05,\,s{=}0.30$ & Yes & 100 & 0.0292 & 0.01605 \\
    Salt-Pepper & $p{=}0.10,\,s{=}0.30$ & No  & 25  & 0.1012 & 0.00345 \\
    Salt-Pepper & $p{=}0.10,\,s{=}0.30$ & Yes & 75  & 0.1104 & 0.02381 \\
    Salt-Pepper & $p{=}0.15,\,s{=}0.30$ & No  & 50  & 0.0768 & 0.00410 \\
    Salt-Pepper & $p{=}0.15,\,s{=}0.30$ & Yes & 50  & 0.0914 & 0.03143 \\
    \bottomrule
  \end{tabular}
\end{table*}

\paragraph{Slip noise.}
The aware policy consistently outperforms naive at all slip levels.
The gap is most dramatic at $\alpha_{\mathrm{slip}} = 0.20$: the naive policy achieves
0\% while the aware policy achieves 100\%.
The mechanism is speed compensation: the 20\% undershoot means the aware policy must
command $v^* / (1 - 0.20) = 1.25\,v^*$ to hit a target at distance $d$.
The naive policy targets $v^*$ and consistently undershoots by 20\% of the required
range.

The non-monotone result at $\alpha_{\mathrm{slip}} = 0.15$ (naive achieves 100\%,
sandwiched between 25\% at $\alpha = 0.10$ and 0\% at $\alpha = 0.20$) is suspicious
and warrants caution.
All results are from a single seed; this level is likely a case of single-seed luck
where the five exploration throws happened to produce a GP with sufficient coverage
of the relevant speed range to allow the naive policy to partially self-correct.
Multi-seed evaluation ($\geq 5$ seeds) is required to confirm whether naive
achieves 100\% at $\alpha = 0.15$ in expectation or whether this is an outlier.

The naive policy's final cost ($\approx 0.001$--$0.004$) is systematically lower
than the aware policy's ($\approx 0.007$--$0.011$).
This is the \emph{false-optimism} consequence of Theorem~\ref{thm:zeromean}: the naive
optimiser evaluates noiseless rollouts and concludes it has solved the task precisely,
while the aware policy honestly reports the expected cost under the noise distribution.
Lower training cost, far worse test performance.

\paragraph{Salt-and-pepper noise.}
The aware policy consistently outperforms naive at $p = 0.05$ and $p = 0.10$.
At $p = 0.15$ both policies plateau at 50\%: the spike probability is high enough that
neither policy can reliably avoid corrupted throws.
Salt-and-pepper noise is partially learnable — the aware policy can reduce the throw
speed to avoid the worst-case spikes — but the irreducible fraction of corrupted throws
sets a performance ceiling.

\paragraph{The 50\% ceiling.}
For the full-slip case ($\alpha_{\mathrm{slip}} = 0.20$, $\sigma = 0.04$\,m/s), the
aware policy plateaus at 50\% (evaluated per-throw, counting each of the 10 evaluation
throws individually as a hit or miss) rather than converging to 100\%.
At a release speed of $\approx 2.5$\,m/s, $\sigma = 0.04$\,m/s produces a landing
standard deviation of roughly 2--4\,cm, which accounts for half the throws missing
the 10\,cm-radius target even with perfect mean aim.
This is irreducible aleatoric variance, not a learning failure.
